\section{Réutilisation du bloc Sync\_Ram}
\label{sec:quadram}

Nous avons ensuite construit \texttt{quadram}, un bloc qui réunit quatre instances de
\texttt{Sync\_Ram} derrière un décodeur. Les RAM partagent les entrées et le bus de sortie ;
le décodeur doit donc sélectionner une seule instance à partir de l'adresse globale.

\subsection{Décodeur d'adresses}

Pour quatre RAM, deux bits de sélection suffisent. Nous avons utilisé une affectation
\texttt{with ... select} pour produire le code à un seul bit actif du
tableau~\ref{tab:decodeur}. Le décodeur est combinatoire : il n'a ni horloge ni état interne.

\begin{table}[H]
\centering
\begin{tabular}{cll}
\toprule
Entrée \texttt{data\_in} & Sortie \texttt{cs\_out(3 downto 0)} & RAM sélectionnée \\
\midrule
\texttt{00} & \texttt{0001} & RAM 0 \\
\texttt{01} & \texttt{0010} & RAM 1 \\
\texttt{10} & \texttt{0100} & RAM 2 \\
\texttt{11} & \texttt{1000} & RAM 3 \\
autres valeurs & \texttt{0000} & aucune \\
\bottomrule
\end{tabular}
\caption{Décodage des deux bits de sélection}
\label{tab:decodeur}
\end{table}

Le cas \texttt{when others} est nécessaire car les bits peuvent notamment prendre les
valeurs indéfinies \texttt{U} ou \texttt{X}. Une adresse indéterminée ne doit pas activer
une RAM. Le testbench applique successivement \texttt{00}, \texttt{01}, \texttt{10},
\texttt{11}, puis \texttt{XX}, pendant 10~ns chacun, et vérifie les sorties par assertions.
Le chronogramme (figure~\ref{fig:decodeur}) reproduit le tableau ; après 40~ns, l'entrée
\texttt{XX} donne bien \texttt{0000}. Les vérifications se terminent à 50~ns.

\simcapture{p3-decoder.png}{660 644 10 130}{660 154 10 630}
{Décodeur 2 vers 4 : une sélection par adresse valide et aucune sélection pour XX}
{fig:decodeur}

\subsection{Assemblage de quatre RAM}

Le fichier \texttt{quadram.vhd} instancie le décodeur et quatre exemplaires du modèle
\texttt{sync\_ram} avec sélection. Chaque RAM est configurée avec
\texttt{Nbits\_addr = 8} et \texttt{Nbits\_data = 8}, soit 256 mots de 8 bits. L'ensemble
contient donc $4\times256=1024$ mots de 8 bits.

L'adresse globale comporte 10 bits. Les bits 9 et 8 sont appliqués au décodeur ; les
bits 7 à 0 désignent le mot à l'intérieur de la RAM. L'adresse locale, les données à
écrire, \texttt{we} et l'horloge sont communes aux quatre instances. Seul le signal
\texttt{cs(i)} est propre à chaque RAM. Le tableau~\ref{tab:quadram-adresses} donne la
répartition de l'espace d'adressage.

\begin{table}[H]
\centering
\begin{tabular}{cll}
\toprule
RAM & Bits 9--8 & Adresses globales (hexadécimal) \\
\midrule
0 & \texttt{00} & \texttt{000} à \texttt{0FF} \\
1 & \texttt{01} & \texttt{100} à \texttt{1FF} \\
2 & \texttt{10} & \texttt{200} à \texttt{2FF} \\
3 & \texttt{11} & \texttt{300} à \texttt{3FF} \\
\bottomrule
\end{tabular}
\caption{Répartition des adresses entre les quatre RAM de 256 mots}
\label{tab:quadram-adresses}
\end{table}

Les sorties sont reliées au même signal \texttt{dataout}. Dans la simulation VHDL, les
éléments résolus de type \texttt{std\_logic} qui composent les vecteurs \texttt{unsigned}
permettent cette connexion à plusieurs pilotes. La RAM sélectionnée fournit la donnée,
les autres fournissent \texttt{Z}. Deux pilotes imposant des valeurs opposées produiraient
une valeur \texttt{X}, d'où la nécessité de n'activer qu'une seule sortie.

\subsection{Validation de l'ensemble}

Nous avons compilé la RAM corrigée, le décodeur et \texttt{quadram} dans une même
bibliothèque, puis exécuté \texttt{quadram\_tb.vhd}. Le test écrit quatre valeurs différentes
au déplacement local 5, puis relit les quatre adresses (tableau~\ref{tab:quadram-test}).

\begin{table}[H]
\centering
\begin{tabular}{clll}
\toprule
RAM & Adresse globale & Donnée écrite et relue & Sélection \\
\midrule
0 & \texttt{005} & \texttt{11} & \texttt{0001} \\
1 & \texttt{105} & \texttt{22} & \texttt{0010} \\
2 & \texttt{205} & \texttt{33} & \texttt{0100} \\
3 & \texttt{305} & \texttt{44} & \texttt{1000} \\
\bottomrule
\end{tabular}
\caption{Test des quatre RAM ; adresses et données en hexadécimal}
\label{tab:quadram-test}
\end{table}

La figure~\ref{fig:quadram-test} montre les écritures de 10 à 50~ns puis les lectures de
50 à 90~ns. Les assertions ne signalent aucune erreur et le message
\og Fin du test de quadram\fg{} apparaît à 90~ns ; l'horloge s'arrête à 100~ns. Les quatre
valeurs restent distinctes : une écriture dans une RAM n'a pas écrasé le mot de même
déplacement dans les autres. Si l'écriture n'était pas conditionnée par \texttt{cs}, les
quatre mots finiraient tous avec la dernière valeur, \texttt{44}, et les trois premières
comparaisons échoueraient.

\simcapture{p3-quadram.png}{660 674 10 130}{660 339 10 550}
{Quadram fixe : écritures puis lectures indépendantes des quatre RAM}
{fig:quadram-test}

\paragraph{Interprétation du temps de lecture.}
Dans ce test, le déplacement local reste égal à 5 pendant le passage d'une RAM à l'autre.
Le changement des bits de sélection modifie donc la sortie dès que le décodeur change
\texttt{cs} : le mot local 5 est déjà adressé dans chaque RAM. Ce comportement ne signifie
pas que toute lecture est asynchrone. Si le déplacement change aussi, la sortie peut
transitoirement présenter la nouvelle RAM à l'ancienne adresse locale, jusqu'au front
montant suivant. Le processeur devra attendre que cette adresse soit échantillonnée avant
de capturer la donnée (section~\ref{sec:synchro}).

Ces quatre lectures vérifient la sélection et l'indépendance des blocs au déplacement 5 ;
elles ne constituent pas un parcours de toutes les adresses de \texttt{quadram}.

\subsection{Généricité du décodeur et de quadram}

Nous avons rendu le décodeur générique avec \texttt{Nbits\_sel}, le nombre de bits de
sélection. Le décodeur produit $2^{\mathtt{Nbits\_sel}}$ sorties, chacune décrite par une
affectation générée (code~\ref{lst:decodeur-generique}).

\noindent\begin{minipage}{\linewidth}
\begin{lstlisting}[style=vhdl, caption={Décodage générique des signaux de sélection}, label={lst:decodeur-generique}]
GEN_CS: for i in 0 to 2**Nbits_sel-1 generate
  cs_out(i) <= '1' when data_in = i else '0';
end generate GEN_CS;
\end{lstlisting}
\end{minipage}\par

Le bloc \texttt{quadram} reçoit trois paramètres : \texttt{Nbits\_sel},
\texttt{Nbits\_addr} et \texttt{Nbits\_data}, notés ici $s$, $a$ et $d$.
Il contient $N=2^s$ RAM de $2^a$ mots de $d$ bits, soit $2^{s+a}$ mots au total ;
l'adresse globale a $s+a$ bits. Le nombre de RAM est donc une puissance de deux :
cette variante ne couvre pas un nombre arbitraire de blocs.

Les instanciations explicites sont remplacées par le code~\ref{lst:quadram-generate}.
\texttt{generate} crée des instances matérielles parallèles à l'élaboration du modèle,
et non une boucle d'exécution partageant une seule RAM.

\noindent\begin{minipage}{\linewidth}
\begin{lstlisting}[style=vhdl, caption={Instances de RAM créées par for ... generate}, label={lst:quadram-generate}]
RAMS: for i in 0 to 2**Nbits_sel-1 generate
  RAM: entity work.sync_ram(RTL)
    generic map (Nbits_addr => Nbits_addr, Nbits_data => Nbits_data)
    port map (cs => cs(i), clock => clock, we => we,
              address => address(Nbits_addr-1 downto 0),
              datain => datain, dataout => dataout);
end generate RAMS;
\end{lstlisting}
\end{minipage}\par

\paragraph{Validation de deux configurations.}
Le testbench écrit au déplacement 5 de la RAM $i$ la valeur \texttt{0x11} multipliée par $i+1$,
puis vérifie la lecture par assertion. Le tableau~\ref{tab:quadram-configs} présente
les deux configurations testées sans modification des sources.

\begin{table}[H]
\centering
\begin{tabular}{rrrrrl}
\toprule
\makecell{Bits de\\sélection} & \makecell{Bits d'adresse\\locale} &
\makecell{Bits de\\données} & RAM & \makecell{Mots\\par RAM} & Résultat \\
\midrule
3 & 4 & 16 & 8 & 16 & aucune assertion en erreur \\
2 & 10 & 32 & 4 & 1024 & aucune assertion en erreur \\
\bottomrule
\end{tabular}
\caption{Configurations du quadram générique effectivement simulées}
\label{tab:quadram-configs}
\end{table}

Avec huit RAM, les adresses \texttt{05}, \texttt{15}, jusqu'à \texttt{75} relisent
\texttt{0011}, \texttt{0022}, jusqu'à \texttt{0088}. Le test se termine à 170~ns et
l'horloge s'arrête à 180~ns (figure~\ref{fig:quadram-g8}).

\simcapture{p3-generic-8x16x16.png}{660 674 10 130}{660 339 10 550}
{Version générique : huit RAM de 16 mots de 16 bits, test des huit sélections}
{fig:quadram-g8}

La configuration du système final contient quatre RAM de 1024 mots de 32 bits.
L'adresse de 12 bits comprend deux bits de sélection et dix bits d'adresse locale.
Les mots \texttt{005}, \texttt{405}, \texttt{805} et \texttt{C05} relisent respectivement
\texttt{11}, \texttt{22}, \texttt{33} et \texttt{44} sur 32 bits. Les assertions passent,
la simulation s'arrête à 100~ns (figure~\ref{fig:quadram-g4}). La sélection et l'isolation
sont ainsi vérifiées au déplacement utilisé dans les deux configurations.

\simcapture{p3-generic-4x1024x32.png}{660 674 10 130}{660 194 10 690}
{Configuration du système final : quatre RAM de 1024 mots de 32 bits}
{fig:quadram-g4}
